\section{负熵远征}

下一代机器可以让答案更容易获得。文明仍要决定维持什么、发现什么，以及准备让谁承担什么代价。

“负熵”首先提醒一项物理事实：生命与工程系统可以使用资源、输出熵流，维持局部组织。它在这里也承载一层人的愿望：把新增能力用于值得持续的世界。物理账本说明维持的条件，不能替我们选择值得维持的状态。

\subsection{形成与维持}

复制给出了一个小而明确的例子。若$\varepsilon_0$为初始每位错误率，校对与修复的条件性残余因子分别为$f_{\text{校对}}$和$f_{\text{修复}}$，在可以采用乘法组合的模型中，
\begin{equation}
\varepsilon_{\text{总}}=\varepsilon_0f_{\text{校对}}f_{\text{修复}},
\qquad
\mathbb E[M]=\varepsilon_{\text{总}}L .
\end{equation}
$M$为长度$L$上的错误数。期望错误数不等于至少发生一次错误的概率；保真度还依赖位置、上下文、相关性与修复机制。\cite{kunkel2000dna}

关键问题是：什么过程让结构在反复调用中保持可用。生物校对、程序测试、档案与机构复核采用不同机制，各自要核算成本和效果。

麦克斯韦妖把信息与物理边界之间的关系推得很深。信息可以引导操作，但循环的完整账本包括测量、存储与复位。兰道尔界限约束相应条件下的信息擦除，Bennett则梳理了计算、可逆性与妖的热力学。\cite{landauer1961,bennett1982thermodynamics} 把存储与热库放在图外，一次聪明操作就可能显得免费。

同样的遗漏也会发生在AI接口。候选便宜，可能因为准备已被摊销、候选可以撤回，或确认发生在别处。这些都是真实技术优势。最终价值要看它们使怎样的完整过程成为可能。

\subsection{丰裕扩大选择}

文明可以用智能改善健康、基础设施、教育、发现与生态，也可以用它加剧注意力竞争，或把某个狭窄目标优化得极其充分。去向由产权、激励、公共安排、组织与个人项目共同塑造。

舒适在这幅图里有正当位置。减轻疼痛、重复劳动、排斥与无用等待，都是进步。生活还需要关系、好奇、自主投入，以及判断能够得到回应的地方。

高阻力前沿提供一部分这样的相遇。一次艰难实验，一套被持续维护的公共服务，一片逐步恢复的生态，一个人的成长，一段关系，一门学到足以判断工具的手艺，也提供同样的入口。规模不同，要求相通：结果要通过行动和后果建立。

\subsection{在约束中扩大自由}

自由需要可用的可能性，对相关边界的理解，以及足以作选择的权力和支持。好的工程降低资源障碍、改善过程、隔离依赖、提高测量并建立恢复路径；好的制度处理新增收益和剩余代价怎样分配。

智能的规模让这项工作更加紧迫。候选便宜以后，选择的质量更重要；行动加快以后，反馈的匹配更重要；作用范围扩大以后，授权和责任更重要。

负熵远征，就是把新增可能变成值得持续的世界。它需要强大的搜索，需要能够行动的身体与机器，需要足以打断说法的证据，也需要愿意承担长期账本的人和机构。

智能寻找路径。文明由它建成的路径、持续看见的后果，以及这些路径使之成为可能的生活塑造。

\furtherreading{%
薛定谔《生命是什么》（Cambridge University Press, 1944）\\
冯·诺依曼《自复制自动机理论》（Burks 编，University of Illinois Press, 1966）\\
Nicholas Negroponte《数字化生存》（Alfred A. Knopf, 1995）\\
William F. Ogburn《社会变迁》（B. W. Huebsch, 1922）\\
William J. Baumol and William G. Bowen《表演艺术：经济的困境》（Twentieth Century Fund, 1966）\\
兰道尔《不可逆性与计算过程中的热产生》（IBM Journal of Research and Development, 1961）\\
Markus Reiner《德博拉数》（Physics Today, 1964）\\
香农《通信的数学理论》（Bell System Technical Journal, 1948）}


\appendix
%% ===================== 本章参考文献 =====================
\printbibliography[heading=subbibliography,title=参考文献]
